\documentclass[11pt]{article}
\usepackage[margin=1in]{geometry}
\usepackage{amsmath}
\usepackage{hyperref}
\title{README: Silicon Radiation Damage Project}
\begin{document}
\maketitle

This project studies the chain
\[
\text{radiation} \rightarrow \text{defects in silicon} \rightarrow \text{material degradation} \rightarrow \text{device degradation}.
\]

The repository is divided into:
\begin{itemize}
    \item Geant4 radiation transport module
    \item MATLAB defect and device degradation module
    \item notes and derivations for presentation preparation
\end{itemize}

The intended workflow is:
\begin{enumerate}
    \item run Geant4 to produce radiation-interaction outputs,
    \item use MATLAB to convert those outputs into defect-density and device-metric trends,
    \item use the notes and derivations to support slides and discussion.
\end{enumerate}

\section*{Main physics idea}
Energetic particles can transfer enough energy to displace silicon atoms from lattice sites. The resulting vacancies, interstitials, and defect complexes create trap and recombination states that reduce carrier lifetime and increase leakage current.

\section*{Cryogenic extension}
An optional extension studies how the kinetics of radiation-induced defects change as temperature is lowered from room temperature toward the cryogenic regime. The extension models:
\begin{itemize}
    \item defect migration freeze-out,
    \item trap emission freeze-out,
    \item annealing freeze-out.
\end{itemize}
The corresponding note files are \texttt{docs/physics\_notes/cryogenic\_defect\_evolution.tex} and \texttt{docs/derivations/cryogenic\_kinetics.tex}. The MATLAB entry point for this section is \texttt{matlab/main\_cryogenic\_extension.m}.

\end{document}
